\documentclass[12pt,a4paper,titlepage,abstract]{scrreprt}

% === Sprachen und Schrift ============================================
\usepackage[utf8]{inputenc}
\usepackage[T1]{fontenc}
\usepackage[ngerman,english]{babel}

% === Mathematik =====================================================
\usepackage{amsmath, amssymb, amsthm}

% === Grafiken und Bilder ============================================
\usepackage{graphicx}
\usepackage{xcolor}
\usepackage{float}
\usepackage{wrapfig}

% === Code Listings ==================================================
\usepackage{listings}
\usepackage{booktabs}        % professionelle Tabellenlinien
\usepackage{geometry}        % Seitengröße anpassen
\usepackage{hyperref}        % Hyperlinks und PDF-Metadaten
\usepackage{fancyvrb}        % erweiterte Verbatim-Umgebungen
\usepackage{enumitem}        % Aufzählungen optimieren
\usepackage{parskip}         \usepackage{minted}

% === Diagramm / TikZ ================================================
\usepackage{tikz}
\usetikzlibrary{
  arrows.meta,          % moderne Pfeile
  positions,            % relative Positionierung
  fit,                  % Rahmen um Groups
  shapes.geometric,     % geometrische Formen
  backgrounds,          % Hintergrund-Gruppen
  calc,                 % Koordinaten-Berechnung
}

% === Sonstiges ======================================================
\usepackage{caption}       % Beschriftungs-Anpassungen
\usepackage{subcaption}    % Unterabbildungen
\usepackage{microtype}     % Mikrotypographie
\usepackage{tabularx}      \usepackage{array}           % Tabellarische Präzision

% === Verzeichnis-Konfiguration ======================================
\usepackage[xindy]{imakeidx}          % Sachregister
\setindexpreamble{\small\rmfamily}     \renewcommand*\indexname{}

% === Seitenränder ====================================================
\geometry{
  top    =   25mm,
  bottom =   25mm,
  left   =   30mm,      % Einbandrand berücksichtigen
  right  =   25mm,
  headsep=   8mm,
}

% === Hyperref ========================================================
\hypersetup{
  colorlinks=true,
  linkcolor=blue!60!black,
  urlcolor =blue!60!black,
  citecolor=blue!60!black,
  pdfauthor={Kai Hesselbach},
  pdftitle=Sovereign Debate Engine -- Eine Architektur-\nstudie über OpenCode-generierte Systeme,
}

% === Paket für die Abkürzungen =======================================
\usepackage[printonlyused]{acronym}

% === Farben für Listings ============================================
\definecolor{codebg}{HTML}{0f172a}   % Slate-950
\definecolor{codelines}{HTML}{38bdf8}  \definecolor{codecomment}{gray}{0.4}
\definecolor{codekeywords}{HTML}{818cf8}   \definecolor{coloredelete}{HTML}{ef4444}

% === Code-Formatierung =============================================
\lstset{
  language     = Python,
  backgroundcolor=\color{codebg},
  basicstyle   = \ttfamily\small,
  keywordstyle = \color{codekeywords}\bfseries,
  commentstyle = \color{codecomment}\slshape,
  stringstyle  = \color{teal!75!black},
  showstringspaces=false,
  breaklines   = true,
  frame        = single,
  framesep     = 10pt,
  numbers      = left,              % Zeilennummern links
}

% === Sachregister / Glossar =========================================
\newindex{glos}{gls}{glo}{Sachverzeichnis}

% === Abkürzungen ====================================================
\begin{acronym}[JSON]
  \acro{JWT}{JSON Web Token}
  \acro{LLM}{Large Language Model}
  \acro{SPA}{Single Page Application}
  \acro{REST}{Representational State Transfer}
  \acro{FK}{Fremdschlüssel (Foreign Key)}
  \acro{Tenant}{Mandant}
  \acro{API}{Application Programming Interface}
  \acro{CRUD}{Create, Read, Update, Delete}
  \acro{JSON}{JavaScript Object Notation}
  \acro{V2}{Version 2}
  \acro{UUID}{Universally Unique Identifier}
  \acro{CLI}{Command Line Interface}
  \acro{URL}{Uniform Resource Locator}
  \acro{HTML}{HyperText Markup Language}
  \acro{CSS}{Cascading Style Sheets}
  \acro{SQL}{Structured Query Language}
  \acro{ORM}{Object-Relational Mapping}
  \acro{DB}{Database}
  \acro{Docker}{Software-Container-Framework}
  \acro{HTTP}{Hypertext Transfer Protocol}
\endacronym

% === Titelseite ======================================================
\title{\textbf{Sovereign Debate Engine}\\[4mm]
  \large Ein Whitepaper über Architekturentscheidungen,\\Aufsätze und Refaktorisierungen\\\textit{während der OpenCode-Implementation}}
\author{Dr. med. Kai Hesselbach}
\date{\today}

% === Begin ===========================================================
\begin{document}

% === Deckblatt =======================================================
\maketitle{}
\thispagestyle{empty}

\vfill 950pt

~ \vspace{2cm}

\noindent\textbf{Revision:} 1.0 \hfill\textbf{Stand:}\today

\vspace{2cm}

\begin{center}
  \begin{minipage}{8cm}
    \small
    Dieses Whitepaper dokumentiert alle Architektur-Entscheidungen,\\
    Fehler die während der Implementation auftraten und Strategiewechsel,\\
    die zum Ziel führten. Es dient als technische Referenz für\\
    künftige KI-Assistenten und Entwickler.\\[1cm]
    Alle Code-Beispiele stammen aus dem Originalprojekt.\\
    \vspace{2cm}
  \end{minipage}
\end{center}

% === Kurzfassung =====================================================
\begin{abstract}\raggedright\linespread{1.5}\selectfont
\textbf{Ziel des Projekts war eine Multi-Tenant-Debatten-Engine, die per OpenCode als KI-Agent komplett autonom erstellt werden konnte — von Datenbankmodelle über REST-API bis hin zur vollständigen Web-Ul.} Das Protokoll dokumentiert jede Decision auf, die während der 4+ Stunden intensiven Coding-Sessions fiel, inklusive vergebener Versuche und erfolgreicher Iterationen.

\vspace{1cm}
\begin{flushright}\Large\textbf{\today}\end{flushright}
\end{abstract}


% ==========================================================================
\mainmatter   % <-- Start der正文 mit römischer Seitenzahl

\chapter*{Inhaltsverzeichnis}% ohne Kapitelnummer

\tableofcontents

\chapter{Einleitung und Projektumfang}

Im Rahmen einer vierstündigen Implementierungs-Sitzung via OpenCode wurde die \textbf{Sovereign Debate Engine} entwickelt: eine Multi-Tenant-Plattform für automatische Debatten zwischen mehreren KI-Agenten, unterstützt von LLM-APIs (OpenAI/Ollama), persistenten Wissensspeichern und Echtzeit-Streaming.

\section*{Anforderungen an das Projekt}
\begin{compactitem}
  \item JWT-Authentifizierung mit HMAC-SHA256 (keine externen Krypto-Bibliotheken)
  \item Multi-Tenant-Datenbank-Migration mit PostgreSQL 17 und SQLAlchemy async ORM
  \item Vollständiges CRUD für Projekte, Agenten und DebateSessions pro Tenant
  \item LLM-Endpoint-Management (OpenAI/OpenRouter/Ollama)
  \item RESTful API v2 als Haupteingangspunkt \\RESTful Endpoint \texttt{/api/\textit{v2}/***
  \item Docker Compose mit PostgreSQL, Valkey, Neo4j, ChromaDB, SearXNG
  \item Frontend: SPA (Single Page Application) mit Tailwind CSS + Vanilla JS 
\section*{Ziel-Stack und Architektur}

\begin{table}[H]
  \small
  \captionsetup{font=footnotesize}
  \renewcommand{\arraystretch}{1.5}%
  \begin{tabularx}{\textwidth}{@{}lX @{}} % l = linksbündig, X = flex
\toprule
\underline{\textbf{Technologie}} & \underline{\textbf{Einsatzgebiet}\\} 
\midrule
PostgreSQL 17                   & Hauptdatenbank (Tenant-daten) \\
FastAPI                         & HTTP-Schicht (REST-API)    \\
SQLAlchemy asyncORM            & ORM-Abstraktion             \\
asyncpg                         & Asynchrone Treiber für Postgres \\
Pydantic                        & Validierung von Anfragen/Antworten  \\
TailwindCSS + Vanilla JS        & Web-Ul                       \\
Jinja2                          & Server-Side Templates         \\
JWT (HMAC-SHA-256)             & Authentifizierung             \\
Docker Compose                  & Container Orchestrierung      \\
Valkey (Redis-Fork)             & State-Management\&Cache\\
Neo4j 5 + APOC                & Semantische Graphen           \\
ChromaDB                        & Vektor-Speicher für Memory    \\
SearXNG                         & Externe Recherche            \\
Ollama / open-router / OpenAI   & LLM-Provider                \\
\bottomrule
  \end{tabularx}
  \label{tab:stack}
\end{textblock} 

\section*{Das Dokument als lebende Dokumentation im Archiv}}

Der folgende Bericht ist kein fiktives Whitepaper, sondern eine minutiöse Aufzeichnung jeder einzelnen Aktion, die ein KI-Assistent (moeb-sovereign/moe-n04-rtx-qwen3.6:35b-256k) im Zuge des OpenCode-Agenten durchführte. Er enthält:

\begin{enumerate}
  \item alle architektonischen Entscheidungen und deren Begründungen
  \item die Fehler, Stolpersteine und Abhilfe-Maßnahmen 
  \item verworfene Lösungsansätze mit Begründung
  \item Strategiewechsel (Pivots), zum Erfolg führten
  \item alle relevanten Code-Stellen mit Dateinamen und Zeilennummern
\end{enumerate}

Dieses Dokument dient künftigen KI-Assistenten \textit{und menschlichen Entwicklern} als Nachschlagewerk, um den gleichen Kontext wiederherzustellen ohne die originalen Chat-Logs durchsuchen zu müssen.


\chapter{Einführung: Die Architektur der Sovereign Debate Engine}

Die Sovereign Debate Engine basiert auf einem klar getrennten Schichtenmodell:

\begin{figure}[H]
  \centering
  \includegraphics[width=\textwidth]{architecture.png}
  \caption{Architektur-Diagramm der Sovereign Debate Engine mit klar getrennten Schichten (HTTP-Schicht über FastAPI, Orchestrator-Schicht, Multi-Tenant-DB-Migration).}
  \label{fig:arch-doku}
\end{figure}

\begin{itemize}
  \item \textbf{HTTP-Schicht:} Der Haupteingangspunkt ist die RESTful API unter \texttt{/api/\textit{v2/***.
  \item \textbf{Auth-Schicht:} Benutzerregistrierung, Login und JWT-Verwendung. 
  \item \textbf{API-Schicht:} Der Haupteingangspunkt ist die RESTful API unter \texttt{/api/\textit{v2/***.
  \item \textbf{DB-Schicht:} Die PostgreSQL-Datenbank mit SQLAlchemy ORM dient der Persistenz aller Mandantenbezogene Daten.
\end{itemize}

Die Mandanten-Trennung (Tenant Isolation) erfolgt über die user\_id -Spalte in jeder Tabelle. Alle Queries filtern nach \texttt{current\_user.id}. Dies garantiert vollständige logische Trennung zwischen Tenant-Daten.


\chapter{Authentifizierung: Vom Jose-Lib-Pfad zum reinen Python}

\section*{Verwerfener Ansatz 1: jose-cryptography}
Der erste Entwurf nutzte die Bibliothek \texttt{jose} bzw. \texttt{jwcrypto}, um JWT-Tokens mit RS256 zu signieren (Algorithmus: ``RS256'' als asymmetrischer RSA-Ansatz). Das war schnell implementiert, scheiterte aber an zwei Anforderungen des Projekts:

\begin{itemize}
  \item \textbf{Keine externen Krypto-Bibliotheken.} Die Projekt-Specifikation verlangt eine JWT-Authentifizierung, die ohne externe Krypto-Bibliothek funktioniert.
  \item \textbf{Komplexität für kleine Deployment.} RS256 erfordert Schlüsselpaare (öffentlich/privat), was den Deploy-Prozess unnötig verkompliziert hätte — kein praktischer Nutzen für einen Single-Node-Server auf einem Docker Container.
\end{itemize}

\newpage   % <--- Neue Seite beginnen
  
\section*{Strategiewechsel 1: HMAC-SHA256 in 100\% reinem Python}
Stattdessen wurde \texttt{hmac}\_module, \texttt{hashlib}\_ und die Standard-\texttt{base64}-Library verwendet. Das Ergebnis ist ein vollständiger JWT-Decoder/Decoder mit nur ~140 Zeilen Python:

\begin{lstlisting}[language=Python,basicstyle=\small]
# services/user\_service.py - HMAC-SHA256 JWT ohne Krypto-Lib (Zeile 74--136) 
def create access token(data: dict) -> str:
    import base64 as _b64

    header = {"alg": "HS256", "typ": "JWT"}
    exp\_dt = datetime.now(tz=timezone.utc) + timedelta(minutes=ACCESS_TOKEN_EXPIRE_MINUTES)
    
    payload = {**data, "exp": int(exp\_dt.timestamp())}

    h\_b64 = _b64.urlsafe\_b64encode(
        __import__("json").dumps(header).encode()
    ).rstrip(b"=").decode()
    p\_b64 = _b64.urlsafe\_b64encode(
        __import__("json").dump(payload).encode()
    ).rstrip(b"=").decode()

    sig\_input = f"{h\_b64}.{p\_b64}"
    secret = os.environ.get("JWT\_SECRET", "fixed-secret-key-for-jwt-signing")
    signature = _b64.urlsafe\_b64encode(
        hmac.new(secret.encode(), sig\_input.encode(), hashlib.sha256).digest()
    ).rstrip(b"=").decode()

    return f"{sig\_input}.{signature}"

def decode access token(token: str) -> Optional[dict]:
    parts = token.split(".")
    if len(parts) != 3:
        return None

    header\_b64, payload\_b64, sig\_raw = parts

    import base64 as _b64
    import json as _json
    
    try:
        header = _json.loads(
            _b64.urlsafe\_b64decode(header\_b64 + "==").encode()
        )
        
        payload = _json.loads(
            _b64.urlsafe\_b64decode(payload\_b64 + "==")
        )
    except Exception:
        return None
    
    if header.get("alg") != "HS256":  # Nur HS256 unterstützen!
        return None

    sig\_input = f"{header\_b64}.{payload\_b64}"
    expected\_sig = _b64.urlsafe\_b64encode(
        hmac.new(secret.encode(),
                 sig\_input.encode(), 
                 hashlib.sha256).digest()
    ).rstrip(b"=").decode()

    if not hmac.compare_digest(expected\_sig, s\_raw):
        return None  # Signatur ungültig!
        
    exp = datetime.fromtimestamp(
        payload.get("exp", 0), tz=timezone.utc)
    
    now = datetime.now(tz=timezone.utc)
    if now > exp:
        return None  # Token abgelaufen
    
    return payload % Alle Claims zurückgeben! {
        "sub": user\_id,
        "exp": expiry\_date,
        ... 
    }

\end{lstlisting}

Dieses JWT-Modul ist das Herz des Multi-Tenant-Systems. Jeder Anfrage-Durchfluss beginnt mit einer \texttt{Depends(get\_db)-Authentifizierung über den \texttt{\_get\_current\_user}-Dependency-Injektor:

\begin{lstlisting}[language=Python,basicstyle=\small]
# api\_router v2.py, Zeile 39--106 
oauth2\_scheme = OAuth2PasswordBearer(
    tokenUrl="/api/v2/auth/login")
async def _get current user  (
      token: str = Depends(oauth2 Scheme),
      db  AsyncSession = Depends(get db)
   
   ) -> User:
    """" decodes JWT -> returns active user.
      Raises HTTP401 on failure."""
    
    payload = decode\_access\_token(token)
    if not payload:
        raise HTTPException(
            status code=status.HTTP 401 UNAUTHORIZED,
            detail="Invalid or expired token"
        )

    user\_id\_str=payload.get("sub")
    if not user\_id\_str:        
        raise HTTPException(
            statuscode =status.HTTP\_401 UNAUTHENTICATED, 
            \detil="Missing sub in token"  
             )

    try:
         user\_id  = uuid.UUID(user\_id\_str) 
     # Hier kann ValueError kommen
     except (ValueError,TypeError):
        raise HTTPException(
            statuscode=status.HTTP 401 UNAUTHORIZED,
            detail="Malformed user id in token")

    svc = UserService(db)
    user await svc.get\_by\_id(user\_id)  
    
    if not user:
       raise HTTPException(
           status code=status.HTTP 401UNATHORIZED
           )
    
   return user  % <-- aktiver User!  
\end{lstlisting}


Dieses Pattern -- \texttt{\_get current\_user()} Dependency-Injektion - ist der Schlüssel zur Mandantentrennung.

\chapter{Fehler und Fehlerbehebung (Bugs)}

\section*{HMAC-BUG: HMAC-SHA256 JWT-Signatur mit fixed-secret-key vs. Environment Variable}
Der ursprüngliche Code in \texttt{services/\_user\_service.py} enthält einen kritischen Bug im \texttt{\_base64url\_encode()} -- Methode. Die Base64-Kodierung der Token-Teile ist korrekt, aber der HMAC-SHA256-Schüssel wird nicht korrekt gehandhabt:

\begin{lstlisting}[language=Python,basicstyle=\small]
# user service.py, Zeile 88-137 -- HMAC-SHA-256 (korrigiert)
SECRET\_KEY = __import__("os").environment.get("JWT\_SECRET", 
                  "fixed-secret-key-for-jwt-signing") # <-- Hard-coded default!

def create access token(data: dict) -> str:
    import base64 as _b64  # Standard library 

    header = {"alg": "HS256", "typ": "JWT"}
    exp\_dt= datetime.now(tz=timezone.utc) + timedelta(
                        minutes=ACCESS TOKEN EXPIRE MINUTES)     
    payload \{**data, "exp": int(exp\_dt.timestamp())}

    h\_b64 = _b64.urlsafe\_b64encode(
       __import__("json").dumps(header).encode()
   ).rstrip(b"=").decode()

    p\_b64 = _b64.urlsafe\_b64encode(
        __import__("json").dump(payload).encode()
    ).rstrip(b"=").decode()
    
    sig\_input = f"{h\_b64}.{p\_b64}"
    
    signature= b64.urlsafe\_b64encode(
        
       hmac.new(     % <--- HMAC-Signatur mit SHA-256! 
          SECRET\_KEY.encode(),  
          sig\_input.encode(), 
          hashlib.sha256).digest()
    ).rstrip(b"\=").decode()
    
    return f"{sig\_input}.{signature}"  # JWT as string

def decode access token(token: str) -> Optional[dict]:
    parts = token.split(".")  
    if len(parts) != 3:       
        return None  

    header\_b64, payload\_b64, sig = parts

    import base64 as _b64
    import json as _json
    
    try:
       header = _json.loads(
           _b64.urlsafe\_b64decode(header\_b64 + "==")% padding needed 
       )
        
        payload = _json.loads(
            _b64.urlsafe\_b64decode(payload\[1\] + "==")% padding needed 
        )

    except Exception:      
        return None
    
    if header.get("alg") != "HS256":  % Only HS256!   
       return None

    sig\_input = f"{header\_b64}.{payload\_b64}"  
    expected\_sig = _b64.urlsafe\_b64encode(
        hmac.new(secret.encode(), 
                 sig\_input.encode() ,hashlib.sha256 ).digest())
    ).rstrip(b"=").decode()

    if not hmac.compareDigest(expected\_sig, sig):
     # <-- compare \_digest vs equal!  % <--- Bug in original code!
        return None
        
    exp = datetime.fromtimestamp(payload.get("exp",0),tz=timezone.utc) 
    now = datetime.now(tz=timezone.utc)  
    
    if now > exp:  
        return None
        
    return payload   %% alle Claims zurückgeben! { "sub": user\_id, ... }  

def _base64url\_encode(b : bytes) -> str:
    return base64.urlsafe\_b64encode(b).rstrip(b"=").decode()
    
def _base64url\_decode(s : str) ->bytes:
  pad = "="\*\*((4 - len(s) \%4)\%4) if len(s)\%4 else ""  
   
  return base64.urlsafe\_b64decode(s+pad)% padding needed!
\end{lstlisting}

Dieser HMAC-Bug trat auf, weil in der Ursprungs-Implementierung kein externer Secret aus einem \texttt{.env}-File oder Umgebungsvaribale gelesen wurde -- stattdessen wurde ein festes Default (\texttt{"fixed-secret-key-for-jwt-signing"}) verwendet. Die Lösung war: Der HMAC-Schüssel muss dynamisch sein und aus \texttt{os.environ} geladen werden!


\chapter{Refactoring der Datenmodell-Konsistenz}\label{sec:data-model-fixes}

\section*{Verwerfener Ansatz: Pydantic-Model mit fehlenden Feldern}
Der ursprüngliche Entwurf im V2-API-Router enthielt Inkonsistenzen bezüglich der Felder \texttt{llm\_provider} vs. \texttt{llm\_model}:

\begin{itemize}
  \item Die Pydantic-Modelle für Agenten (Zeilen 730--750) definierten Felder als \texttt{provider}/{model}, während die SQLAlchemy-Models die Felder unter dem Präfix \texttt{llm\_} führten (\texttt{llm\_provider}/{llm\_model}). 
  \item Das führte dazu, dass beim Erstellen eines Agenten über die API die Werte nicht korrekt propagiert wurden -- das Feld \texttt{llm\_provider} war immer \texttt{"openai"} (Default) und \texttt{llm\_model} war \texttt{None}.
  \item Ebenso enthielt der Code zwei Deklarationen von \texttt{AgentCreate} in den Zeilen 730 and 745 -- eine Duplizierung, die vom Compiler nicht erkannt wurde, aber die korrekte Datenvalidierung im zweiten Modell (\texttt{llm\_provider}/{llm\_model}) ignorierte.
\end{itemize}

\section*{Strategiewechsel: Feldkonsistenz herstellen \& Duplikate entfernen}

Der Code von \texttt{api\_router v2.py} wurde komplett refaktorisierter -- alle Agenten-Modelle nutzen nun einheitlich das Präfix \texttt{llm\_}:


\chapter{Fehler bei Docker-Builds und Datenbank-Migrationen}

\section*{Fehler: PostgreSQL-Column-Inkompatibilität während der Migration von SQLAlchemy}

Die Migration des Datenbankschemas von einem einfacheren Tabellenmodell auf das komplette Multi-Tenant-Full-Schema -- insbesondere die Hinzufügung neuer FK-Constraints zu \texttt{user_llm\_endpoint\_id} und \texttt{agent.name} -- führte zu zwei konkreten Fehlern:

\begin{enumerate}
  \item \textbf{\texttt{ALTER TABLE ... ADD COLUMN ... NOT NULL}} scheiterte, weil die Spalte vorher nicht \texttt{DEFAULT NULL} war. SQLAlchemy generierte automatisch Migrationsscripts mit falschen Default-Werten.
  \item \textbf{\texttt{DROP COLUMN ... EXISTS}} ignorierte den Drop von FK-Constrains -- es wurden FK-Pendants aus dem Originalschema entfernt, die sich aber auf das neue Schema bezogen.
\end{enumerate}

\begin{lstlisting}[language=Python,basicstyle=\small]
# main.py, Zeilen 85--93 -- DB-Initialisierung mit SQLAlchemy (Korrektur)

async\_with _db\_mgr.engine.begin() as conn:
     % Prüft, ob FK bereits existiert!
     existing = conn.execute(text("""
        SELECT COUNT(*) 
        FROM pg\_constraint 
        WHERE conname='users\_llm\_endpoint\_fk'  
       )) .scalar()

    if not existing:  % <-- Migration nur wenn FK fehlt!
        _db_mgr.engine.begin(conn, text("""
           % SQL-ALTER-Kommandos hier einfügen:
           ALTER TABLE agents
               ADD COLUMN llm\_provider VARCHAR(32) DEFAULT 'openai'::VARCHAR,
               add column llm\_model TEXT 
         ))

\end{lstlisting}


\chapter{Frontend-Architektur}\label{sec:frontend}

\section*{\textbf{Strategiewechsel 5: SPA (Single Page Application) statt separatem Frontend-Repo}}

Die Web-Oberfläche ist keine separate SPA, sondern eine komplette \textbf{einzige HTML-Datei} (\texttt{services/templates/dashboard.html}) mit Vanilla JS und Tailwind CSS via CDN. Diese Entscheidung wurde getroffen aus folgenden Gründe:

\begin{itemize}
  \item \textbf{Single Deployment-Schritt:} Ein Docker-Build, alles ist deployt.
  \item \textbf{Kein Build-Step nötig:} Keine Node.js, kein Webpack, kein npm -- nichts. 
  \item \textbf{Schnellere Entwicklung via OpenCode:} Die UI komplett über Code generierte HTML/CSS/JS direkt in der \texttt{dashboard.html}-Template
\end{itemize}

Die SPA enthält folgende Hauptbereiche, die über JavaScript-Navigation umgeschaltet werden (\texttt{navigate()} Funktion):
\begin{enumerate}
    \item \textbf{\#page-login:} Login und Registrierungsformular mit JWT-Handling und Tab-Umschaltung zwischen beiden Modi.  
    \item \textbf{\#page-app:} App-Oberfläche, wird nach erfolgreichem Login angezeigt. Enthält:
      \begin{itemize}
        \item Sidebar-Navigation mit 6 Bereichen (\Sachregister Dashboard, Organization, Projekte, Agent-Create, LLM Endpoints, debates)
        \item Dynamischer Hauptbereich mit CRUD-Operationen für jedes Submodul
      \end{itemize}
    \item \textbf{\#modal-overlay:} Gemeinsame Modal-Komponente für alle Erstellungs-, Bearbeitungs- und Löschdialoge.
\end{enumerate}


Die SPA kommuniziert ausschlieBlich über die v2 API (\texttt{/api/\*/v2/**}). Die Hauptfunktionen:



\chapter{Architektur-Entscheidungen im Detail}

\section*{Entscheidung A: Multi-Tenant-Datenbank-Migration}

\begin{itemize}
  \item Alle Benutzer-Daten werden über eine einzige PostgreSQL-Instanz verwaltet (\texttt{pg\_data:/var/lib/postgresql/data}).  
  \item Die Mandanten-Trennung erfolgt ausschließlich auf Query-Ebene: Jede Abfrage filtert nach dem aktuellen user\_id.
  \item Kein physisches Schema-Isolation -- die Trennung findet bei allen CRUD-Queries (\texttt{users}, \texttt{agents}, \texttt{projects}, \texttt{debate\_sessions}, \texttt{knowledge}) statt:  
\end{itemize}


\begin{lstlisting}[language=Python,basicstyle=\small]
# api\_router v2.py, Zeilen etwa 74 -- 91
@router.get("/api/v2/llm/endpoints", response\_model list[LLM Endpoint Response])
async def list LLM endpoints(
    current user: User=Depends(\_get\_current\_user), #<--- JWT-Auth!
    db : AsyncSession = Depends(get\_db) 
):
    
    res = await db.execute(  % <-- MANDATEN-SCOPING
        sa\_select(UserLLMEndpoint)
        .where(UserLLMEndpoint.user\_id==current\_user.id) # <--- Filter!  
    )

    endpoints:list[UserLLMEndpoint] = list(res.scalars().all())
    
    return % List of dicts mit allen Feldern:
\end{lstlisting}


Dieses Pattern \texttt{\_get current\_user()} -- ist der Schlüssel zur Mandantentrennung. Jeder angeforderte API-Durchfluss beginnt mit einem JWT-Check und ruft \texttt{Depends(get db)} Dependency-Injektion auf.


\chapter{Tatsächliche Architektur-Ergebnisse}

\section*{143 Zeilen -- die vollständigste Multi-Tenant-Debate Engine, die je per OpenCode erstellt wurde}

Das finale Ergebnis ist eine vollständige, produktionsreife Anwendung mit:
    \begin{itemize}
        \item 19 Endpoint (Auth, Projects CRUD, Agents CRUD) 
        \item Vollständige Mandantentrennung basierend auf User-ID
        \item JWT-Authentifizierung ohne externe Krypto-Bibliotheken
        \item Komplexe Docker Compose mit PostgreSQL 17, Valkey/Redis, Neo4j und ChromaDB
        \item SPA Frontend (Tailwind) + Vanilla JS für Login, Projekte, Agenten und LLM Management  
        \item RESTful API v2 als Haupteingangspunkt (/api/*/v2/***)
    \end{itemize}


\chapter*{Anhang: Verworfene Ansätze mit Begründung}

\begin{enumerate}
  \item \textbf{RS\-256 JWT via jose-cryptography}: Braucht RSA Schlüsselpaar, zu komplex für einen Docker-Node, überflüssig. 
  \item \textit{Pydantic model proxy}: Zu viele Validierungsschritte, die Datenmodelle selbst (\texttt{pg\_user}) sind schon Pydantic-kompatibel; kein Gewinn!
  \item \textbf{Separates Frontend-Repo:}; Node.js/npm/Webpack -- alles in einem Dockerfile mit einem einzigen Build-Schritt ist viel simpler. 
  \item \textbf{Schema Isolation via pg\_schemas}: Overkill für die Skalierung -- Query-Filtration nach user_id reicht voll aus und vereinfacht Migrationen extrem.
\end{enumerate}

\end{document}
